\documentclass[10pt]{beamer}

\usepackage{verbatim}
%\usepackage[latin1]{inputenc}
\usepackage{amssymb,amsmath,amsfonts,url}
\usepackage{graphicx,color}
\usepackage{pgf,pgfarrows,pgfnodes,pgfautomata,pgfheaps,pgfshade}
\usepackage[utf8]{inputenc}
\usepackage[french]{babel}
\usepackage{tikz}

\usepackage[all]{xy}
\usepackage{amstext}
\usepackage{xspace}
\usepackage{mdwtab}


\mode<article> {
  \usepackage{fullpage}
  \usepackage{hyperref}
  \setjobnamebeamerversion{advancedcourse.beamer}
}





\newcommand{\alertonce}[1]{\addtocounter{beamerpauses}{-1}\alert<+>{#1}}
\newcommand{\alerttwice}[1]{\addtocounter{beamerpauses}{-1}\alert<+-+(1)>{#1}}

%\newcommand{\subsubsubsection}[1]{\noindent\textbf{#1}}

\newcommand{\cfbox}[1]{\begin{center}\fbox{#1}\end{center}}
\newcommand{\cbox}[1]{\begin{center} #1 \end{center}}

%%%%%%%%%%%%%%%%%%%%%%%%%%%%%%%%%%%%%%%%%%%%%%%%%%%%%%%%%%%%%%%%%%%%%%%%%%%%%%%%%%%%%%%%%

\graphicspath{{Graphics/}}
%\input def_0.tex

%\input{StyleSlides.tex}
\mode<presentation> {
    \usepackage{beamerthemeBoadilla}

  %\beamertemplatetransparentcovereddynamic
  \beamertemplateballitem
}
\usepackage{listings}
\usepackage{xcolor}

\lstset{
	language=Python,
	basicstyle=\ttfamily\small,
	keywordstyle=\color{blue},
	commentstyle=\color{gray},
	stringstyle=\color{green!50!black},
	showstringspaces=false,
	frame=single,
	numbers=none
}
\lstset{
	language=Python,
	basicstyle=\ttfamily\small,
	keywordstyle=\color{blue},
	commentstyle=\color{gray},
	stringstyle=\color{green!50!black},
	showstringspaces=false,
	frame=single,
	numbers=none,
	literate=
	{é}{{\'e}}1
	{è}{{\`e}}1
	{ê}{{\^e}}1
	{ë}{{\"e}}1
	{à}{{\`a}}1
	{â}{{\^a}}1
	{ä}{{\"a}}1
	{ù}{{\`u}}1
	{û}{{\^u}}1
	{ü}{{\"u}}1
	{ô}{{\^o}}1
	{ö}{{\"o}}1
	{î}{{\^i}}1
	{ï}{{\"i}}1
	{ç}{{\c c}}1
}



\title[IN100 Cours 9]{\textbf{Fondements informatiques I}\\
Cours 9: Algorithmes : manipuler des tableaux}

\author[]{Sorina Ionica \texttt{sorina.ionica@uvsq.fr} \\~\\Bilal FAYE  \texttt{bilal.faye@uvsq.fr} }


\institute{}

\date{}


%\setbeamercolor{normal text}{fg=black}

\begin{document}


\frame{\titlepage}
\begin{frame}{Qu'est-ce qu'un algorithme ?}

Un \textbf{algorithme} est une suite finie d'étapes permettant de
résoudre un problème ou d'effectuer une tâche.

\begin{exampleblock}{Exemple : analyser les notes d'un examen}
On dispose d'un tableau de notes :
\[
\texttt{notes = [12, 8, 15, 10, 6]}
\]

On souhaite :
\begin{itemize}
    \item trouver la note la plus petite ;
    \item trouver la note la plus grande ;
    \item afficher les notes dans l'ordre croissant.
\end{itemize}
\end{exampleblock}

\pause

\begin{block}{Les trois éléments d'un algorithme}
\begin{itemize}
    \item \textbf{Entrée :} les données dont on dispose (ici, les notes).
    \item \textbf{Traitement :} les étapes à effectuer.
    \item \textbf{Sortie :} le résultat attendu.
\end{itemize}
\end{block}

\end{frame}
\begin{frame}{Avant de coder : réfléchir à la méthode}

Pour résoudre un problème, il faut d'abord réfléchir
aux étapes nécessaires, avant de les traduire en Python.

\begin{exampleblock}{Exemple : trouver la plus petite note}
Comment déterminer la plus petite valeur d'un tableau de notes ?

\begin{enumerate}
    \item Prendre la première note comme minimum provisoire.
    \item Parcourir les autres notes une par une.
    \item Si une note est plus petite, mettre à jour le minimum.
    \item À la fin, le minimum provisoire est la plus petite note.
\end{enumerate}
\end{exampleblock}

\pause

\begin{alertblock}{Une bonne habitude}
Avant de programmer, se poser trois questions :
\begin{itemize}
    \item Quelles sont les données en entrée ?
    \item Quelles étapes permettent de résoudre le problème ?
    \item Quel résultat doit-on obtenir ?
\end{itemize}
\end{alertblock}

\end{frame}
\begin{frame}{Minimum d'un tableau : l'algorithme}

On dispose d'un tableau d'entiers. On souhaite trouver
\textbf{l'indice de la plus petite valeur}.

\begin{exampleblock}{Exemple}
\[
L = [12,\ 8,\ 15,\ 6,\ 10]
\]
La plus petite valeur est $6$, et son indice est $3$.
\end{exampleblock}

\pause

\begin{block}{Algorithme}
\begin{itemize}
    \item \textbf{Entrée :} un tableau $L$ d'entiers non vide.
    \item \textbf{Sortie :} l'indice du plus petit élément de $L$.
\end{itemize}

\medskip
\begin{enumerate}
    \item Initialiser $m$ à $0$ : on suppose que le premier élément est le minimum.
    \item Parcourir les indices $i$ de $1$ à $\text{longueur}(L)-1$.
    \item Si $L[i] < L[m]$, alors mettre $m$ à jour : $m \leftarrow i$.
    \item Renvoyer $m$.
\end{enumerate}
\end{block}

\end{frame}
\begin{frame}[fragile]{Minimum d'un tableau : en Python}

Les étapes de l'algorithme se traduisent directement en Python :

\begin{exampleblock}{Implémentation}
\begin{lstlisting}[language=Python]
def minimum(L):
    m = 0
    for i in range(1, len(L)):
        if L[i] < L[m]:
            m = i
    return m
\end{lstlisting}
\end{exampleblock}
\begin{exampleblock}{Exemple d'utilisation}
\begin{lstlisting}[language=Python]
L = [12, 8, 15, 6, 10]
indice = minimum(L)
print(indice)    # 3
print(L[indice]) # 6
\end{lstlisting}
\end{exampleblock}

\textbf{Question :} comment adapter cet algorithme pour trouver
l'indice du plus grand élément ?

\end{frame}
\begin{frame}{Trier un tableau : le tri par sélection}

On souhaite ranger les éléments d'un tableau dans l'ordre
croissant.

\begin{exampleblock}{Exemple}
\[
T = [5,\ 1,\ -8,\ 2,\ -4,\ 7]
\]
On veut obtenir :
\[
[-8,\ -4,\ 1,\ 2,\ 5,\ 7]
\]
\end{exampleblock}

\pause

\begin{block}{Principe du tri par sélection}
On construit progressivement la partie triée du tableau :

\begin{enumerate}
    \item On cherche le plus petit élément du tableau.
    \item On le place au début, en l'échangeant avec le premier élément.
    \item On recommence avec la partie restante du tableau,
    en ignorant les éléments déjà placés.
\end{enumerate}
\end{block}

\end{frame}
\begin{frame}{Tri par sélection : déroulement}

À chaque étape, on cherche le minimum dans la partie
du tableau qui n'est pas encore triée, puis on l'échange
avec le premier élément de cette partie.

\begin{exampleblock}{Exemple : $T = [5,\ 1,\ -8,\ 2,\ -4,\ 7]$}
\small
\begin{itemize}
    \item Départ : \quad $[5,\ 1,\ -8,\ 2,\ -4,\ 7]$
    \pause
    \item Étape 1 : minimum $=-8$ \quad
    $[-8,\ 1,\ 5,\ 2,\ -4,\ 7]$
    \pause
    \item Étape 2 : minimum $=-4$ \quad
    $[-8,\ -4,\ 5,\ 2,\ 1,\ 7]$
    \pause
    \item Étape 3 : minimum $=1$ \quad
    $[-8,\ -4,\ 1,\ 2,\ 5,\ 7]$
    \pause
    \item Étapes suivantes : le tableau reste inchangé.
\end{itemize}
\end{exampleblock}

\pause

\begin{alertblock}{À retenir}
À chaque étape, un élément est placé définitivement
à sa position correcte. La partie triée grandit progressivement.
\end{alertblock}

\end{frame}
\begin{frame}{Tri par sélection : l'algorithme}

\begin{block}{Algorithme}
\begin{itemize}
    \item \textbf{Entrée :} un tableau $T$.
    \item \textbf{Sortie :} le tableau $T$ trié dans l'ordre croissant.
\end{itemize}

\medskip

Pour chaque indice $i$ de $0$ à $\text{longueur}(T)-2$ :

\begin{enumerate}
    \item Chercher l'indice $m$ du minimum dans $T[i:]$.
    \item Échanger $T[i]$ et $T[i+m]$.
\end{enumerate}
\end{block}

\pause

\begin{alertblock}{Pourquoi $i+m$ ?}
La recherche du minimum se fait dans le sous-tableau $T[i:]$.
L'indice $m$ est relatif à ce sous-tableau : on ajoute donc $i$
pour obtenir l'indice correspondant dans le tableau initial.
\end{alertblock}

\end{frame}


\begin{frame}[fragile]{Tri par sélection : en Python}

On réutilise la fonction \texttt{minimum} vue précédemment.

\begin{exampleblock}{Implémentation}
\begin{lstlisting}[language=Python]
def tri_selection(T):
    for i in range(len(T) - 1):
        m = minimum(T[i:])
        T[i], T[i + m] = T[i + m], T[i]
    return T
\end{lstlisting}
\end{exampleblock}


\begin{exampleblock}{Exemple d'utilisation}
\begin{lstlisting}[language=Python]
T = [5, 1, -8, 2, -4, 7]
print(tri_selection(T))
# [-8, -4, 1, 2, 5, 7]
\end{lstlisting}
\end{exampleblock}

\end{frame}
\begin{frame}{Construire un tableau trié progressivement}

Une autre méthode consiste à construire le tableau trié
\textbf{au fur et à mesure que l'on reçoit les éléments}.

\begin{exampleblock}{Exemple}
On reçoit successivement les nombres :
\[
5,\quad 1,\quad -8,\quad 2,\quad -4,\quad 7
\]
On veut construire un tableau trié dans l'ordre croissant.
\end{exampleblock}

\pause

\begin{block}{Principe : l'insertion}
À chaque nouvel élément :
\begin{enumerate}
    \item On cherche sa place dans le tableau déjà trié.
    \item On décale les éléments suivants vers la droite pour libérer une place.
    \item On insère le nouvel élément à sa bonne position.
\end{enumerate}
\end{block}

\pause

\begin{alertblock}{Idée essentielle}
Le tableau est trié après chaque insertion.
\end{alertblock}

\end{frame}
\begin{frame}{Insertion dans un tableau trié : exemple}

Supposons que le tableau soit déjà trié :
\[
T = [-8,\ 1,\ 2,\ 5]
\]

On souhaite y insérer le nombre $-4$.

\pause

\begin{enumerate}
    \item On parcourt le tableau jusqu'à trouver le premier élément
    supérieur ou égal à $-4$ : ici, $1$.
    \pause
    \item On décale les éléments à partir de cette position vers la droite.
    \pause
    \item On place $-4$ dans la case libérée.
\end{enumerate}

\pause

\begin{exampleblock}{Résultat}
\[
[-8,\ 1,\ 2,\ 5]
\quad\longrightarrow\quad
[-8,\ \boxed{-4},\ 1,\ 2,\ 5]
\]
\end{exampleblock}

\end{frame}
\begin{frame}{Insertion dans un tableau trié : algorithme}

On suppose que le tableau $T$ est déjà trié dans l'ordre croissant.

\begin{block}{Algorithme}
\begin{itemize}
    \item \textbf{Entrée :} un tableau trié $T$ et un élément $t$.
    \item \textbf{Sortie :} le tableau $T$ avec $t$ inséré à la bonne place.
\end{itemize}

\medskip

\begin{enumerate}
    \item Initialiser $i$ à $0$.
    \item Tant que $i < \text{longueur}(T)$ et $T[i] < t$,
    avancer à l'indice suivant.
    \item Décaler vers la droite les éléments à partir de l'indice $i$.
    \item Placer $t$ à l'indice $i$.
\end{enumerate}
\end{block}

\pause

\begin{alertblock}{Comment trouver la position ?}
On s'arrête dès que l'on rencontre un élément supérieur ou égal à $t$,
ou lorsque l'on atteint la fin du tableau.
\end{alertblock}

\end{frame}
\begin{frame}{Construire un tableau trié : exemple complet}

On insère successivement les éléments de la séquence :
$5,\ 1,\ -8,\ 2,\ -4,\ 7$.

\small
\begin{exampleblock}{Évolution du tableau}
\begin{itemize}
    \item Insertion de $5$ : \quad $[5]$
    \pause
    \item Insertion de $1$ : \quad $[1,\ 5]$
    \pause
    \item Insertion de $-8$ : \quad $[-8,\ 1,\ 5]$
    \pause
    \item Insertion de $2$ : \quad $[-8,\ 1,\ 2,\ 5]$
    \pause
    \item Insertion de $-4$ : \quad $[-8,\ -4,\ 1,\ 2,\ 5]$
    \pause
    \item Insertion de $7$ : \quad $[-8,\ -4,\ 1,\ 2,\ 5,\ 7]$
\end{itemize}
\end{exampleblock}

\pause

\begin{alertblock}{À retenir}
Après chaque insertion, le tableau reste trié.
À la fin, tous les éléments ont été insérés.
\end{alertblock}

\end{frame}
\begin{frame}[fragile]{Insertion dans un tableau trié : en Python}

On utilise \texttt{insert(i, t)} pour insérer l'élément
\texttt{t} à l'indice \texttt{i}. Les éléments suivants
sont automatiquement décalés vers la droite.

\begin{exampleblock}{Insérer un élément}
\begin{lstlisting}[language=Python]
def insertion(t, T):
    i = 0
    while i < len(T) and T[i] < t:
        i += 1
    T.insert(i, t)
\end{lstlisting}
\end{exampleblock}

\pause

\begin{exampleblock}{Construire un tableau trié}
\begin{lstlisting}[language=Python]
T = []
for t in [5, 1, -8, 2, -4, 7]:
    insertion(t, T)

print(T)
# [-8, -4, 1, 2, 5, 7]
\end{lstlisting}
\end{exampleblock}

\end{frame}
\begin{frame}{Suite de Fibonacci}

On étudie une suite de nombres dans laquelle chaque terme,
à partir du troisième, est obtenu en additionnant les deux
termes précédents.

\begin{block}{Définition}
La suite de Fibonacci est définie par :
\[
F_0 = 0,\qquad F_1 = 1
\]
et, pour tout $n \geq 2$ :
\[
F_n = F_{n-1} + F_{n-2}
\]
\end{block}

\pause

\begin{exampleblock}{Premiers termes}
\[
\begin{array}{c|cccccccc}
n   & 0 & 1 & 2 & 3 & 4 & 5 & 6 & 7\\
\hline
F_n & 0 & 1 & 1 & 2 & 3 & 5 & 8 & 13
\end{array}
\]
\end{exampleblock}

\pause

\textbf{Objectif :} calculer les termes de Fibonacci jusqu'au rang $n$.

\end{frame}
\begin{frame}{Fibonacci : construire le tableau des termes}

Pour calculer tous les termes jusqu'au rang $n$,
on mémorise les valeurs dans un tableau.

\begin{block}{Principe}
\begin{enumerate}
    \item Initialiser le tableau avec les deux premiers termes :
    $[0,1]$.
    \item Calculer le terme suivant en additionnant les deux
    derniers termes du tableau.
    \item Ajouter ce nouveau terme au tableau.
    \item Recommencer jusqu'au rang $n$.
\end{enumerate}
\end{block}

\pause

\begin{exampleblock}{Exemple : calcul jusqu'au rang 5}
\[
\begin{array}{c|l}
\text{Étape} & \text{Tableau}\\
\hline
\text{Initialisation} & [0,1]\\
\text{Ajout de } F_2 & [0,1,1]\\
\text{Ajout de } F_3 & [0,1,1,2]\\
\text{Ajout de } F_4 & [0,1,1,2,3]\\
\text{Ajout de } F_5 & [0,1,1,2,3,5]
\end{array}
\]
\end{exampleblock}

\end{frame}
\begin{frame}{Fibonacci : calculer seulement le terme $F_n$}

Pour calculer le terme suivant, on a besoin que
des deux termes précédents.

\pause

\begin{block}{Nouvelle idée}
Au lieu de mémoriser tout le tableau, on utilise deux variables :
\begin{itemize}
    \item $a$ : le terme précédent ;
    \item $b$ : le terme courant.
\end{itemize}
À chaque étape, on calcule le terme suivant et on met à jour
les deux variables :
$
(a,b) \leftarrow (b,\ a+b)
$
\end{block}

\pause

\begin{exampleblock}{Exemple : calcul de $F_5$}
\footnotesize
\[
\begin{array}{c|c|c}
\text{Étape} & a & b\\
\hline
\text{Départ} & 0 & 1\\
1 & 1 & 1\\
2 & 1 & 2\\
3 & 2 & 3\\
4 & 3 & 5\\
5 & 5 & 8
\end{array}
\]
À la fin, $a=F_5=5$.
\end{exampleblock}

\end{frame}
\begin{frame}[fragile]{Suite de Fibonacci : en Python}

On utilise deux variables pour calculer directement
le terme demandé.

\begin{exampleblock}{Implémentation}
\begin{lstlisting}[language=Python]
def fibonacci(n):
    a, b = 0, 1
    for i in range(n):
        a, b = b, a + b
    return a
\end{lstlisting}
\end{exampleblock}

\pause

\begin{exampleblock}{Exemples}
\begin{lstlisting}[language=Python]
print(fibonacci(5))  # 5
print(fibonacci(7))  # 13
\end{lstlisting}
\end{exampleblock}

\end{frame}


\begin{frame}{Bilan : concevoir un algorithme efficace}

Deux méthodes permettent de calculer les termes de Fibonacci :

\begin{itemize}
    \item \textbf{Avec un tableau :} on mémorise tous les termes.
    \item \textbf{Avec deux variables :} on ne conserve que les deux
    derniers termes.
\end{itemize}

\pause

\begin{alertblock}{À retenir}
Un bon algorithme doit non seulement donner le bon résultat,
mais aussi utiliser efficacement les ressources :
\begin{itemize}
    \item limiter la mémoire utilisée ;
    \item limiter le nombre d'opérations effectuées.
\end{itemize}
\end{alertblock}

\end{frame}
\begin{comment}
\frame[containsverbatim]{
\frametitle{Qu'est-ce que c'est qu'un algorithme?}

\begin{alertblock}{Rappel (premier cours)} 
Un algorithme est une suite finie d'instructions et opérations permettant de résoudre un problème. 
Pour donner un algorithme: 
 
 \begin{itemize}
 \item Input: les données en entrée du problème
 \item Output : la solution du problème
 \item Les étapes qu'on doit effectuer pour résoudre
 \end{itemize}
 \end{alertblock}

\vfill

Exemple concret :  Dans un tableau de notes des étudiants de L1, on veut calculer la note la plus petite et la plus grande. \\

\vfill

Ensuite on voudrait faire afficher le tableau (la liste) des notes  à un examen,  triée dans l'ordre croissante. 


}


\frame[containsverbatim]{
\frametitle{Pour quoi?}

\textcolor{blue}{La bonne pratique:} avant de commencer à coder, on réfléchit sur papier!\\

\vspace{0.5 cm}

En langage humain, on précise les étapes de notre algorithme (pseudocode). \\

\vspace{0.5 cm}

On réfléchit également à la structure du programme :  \\

\vspace{0.3 cm}
\begin{center}
Quelles variables? De quel type? Quelles fonctions? \\
\end{center}
\vspace{0.5 cm}

Pour chaque fonction on précise les paramètres en entrée et en sortie. \\

}


\frame[containsverbatim]{
\frametitle{Minimum d'un tableau}


On veut calculer le minimum d'un tableau et renvoyer son indice. 


\begin{alertblock}{Algorithme}
\begin{verbatim} 
Input: un tableau L d'entiers
Output: l'indice du plus petit entier dans la liste

Au début on initialise m=0 (premier indice du tableau)
Pour tout i allant de 1 à longueur(L)-1
     Si L[j] < L[m] alors
            m <- j 
On renvoit m
\end{verbatim}
\end{alertblock}

}



\frame[containsverbatim]{
\frametitle{Minimum d'un tableau}


En Python cela donne: 


\begin{exampleblock}{}
\begin{verbatim} 
def minimum(L):
    m = 0
    for j in range(1, len(L)):
        if L[j] < L[m]:
            m = j
    return m
\end{verbatim}
\end{exampleblock}

\vspace{0.5 cm}

Et pour calculer le maximum? 

}




\begin{frame}[fragile]{Tri par séléction}



\begin{exampleblock}{Algorithme}
\begin{verbatim}
Input: le tableau L
Output: le tableau L trié

Pour tout i allant de 0 à longueur(L)-1: 
        On calcule l'indice m du plus petit élément dans L[i:]
        On échange L[i] et L[m]
\end{verbatim}
\end{exampleblock}

\vspace{0.5 cm}

\pause

\small{
Exécution sur le tableau  $L = [5, 1, -8, 2, -4, 7]$ :

\begin{itemize}
\item $i=0$    minimum(L)=2 alors $L = [-8, 1, 5, 2, -4, 7]$ 
\pause
\item $i=1$    minimum(L[1:])=4 alors $L = [-8, -4, 5, 2, 1, 7]$ 
\pause
\item $i=2$    minimum(L[2:])=2 alors  $ L = [-8, -4, 1, 2, 5, 7]$ 
\pause
\item $i=3$    minimum(L[3:])=3 alors  $ L = [-8, -4, 1, 2, 5, 7]$ 
\pause
\item $i=4$    minimum(L[4:])=4 alors  $ L = [-8, -4, 1, 2, 5, 7]$ 
%\item $i=5$    minimum(L[5:])=3 alors  $ L = [-8, -4, 1, 2, 5, 7]$ 
\end{itemize}
}
\end{frame}

\frame[containsverbatim]{
\frametitle{Tri par sélection en Python}


\medskip
\begin{alertblock}{}
\begin{verbatim}
def tri_selection(L):
    for i in range(len(L)):
        m = minimum(L[i:])
        L[i], L[m+i] = L[m+i], L[i] # échanger L[i] et L[m]
    return L

L = [5, 1, -8, 2, -4, 7]
print(tri_selection(L))
\end{verbatim}
\end{alertblock}


}




\frame[containsverbatim]{
\frametitle{Comment créer un tableau trié directement}

Si le tableau est crée au fur et à mesure, on insère les éléments dans le bon ordre.


\begin{alertblock}{Algorithme}
\begin{verbatim}
Input : un tableau Tab trié, un élément t
Output:  le tableau trié Tab modifiée après rajout de t


On initialise i<-0. 
taille<-longueur(Tab)
Tant que i<taille and t>l[i]
       On avance dans la liste i=i+1
On décale les éléments Tab[i], ..., Tab[taille-1] d'un cran     
On place l'élément t à l'indice i dans tab (Tab[i]=t)
On renvoie Tab       
\end{verbatim}

\end{alertblock}



}



\begin{frame}[fragile]{Comment créer un tableau trié directement}



\begin{alertblock}{Algorithme}
\begin{verbatim}
On initialise i=0. 
taille<-longueur(Tab)
Tant que i<taille and t>l[i]
      On avance dans la liste i=i+1
On décale les éléments Tab[i], ..., Tab[taille-1] d'un cran     
On place l'élément t à l'indice i dans Tab (Tab[i]=t)
On renvoie Tab   
\end{verbatim}

\end{alertblock}

\small{
Exécution pour construire  un tableau trié à partie de la séquence  5, 1, -8, 2, -4, 7:

\pause
\begin{itemize}
\item $i=0$  alors    $Tab=[5]$
\pause
\item $i=1$ alors     $Tab=[1,5]$
\pause
\item $i=2$ alors   $Tab=[-8,1,5]$
\pause
\item $i=3$ alors    $Tab=[-8,1,2,5]$
\pause
\item $i=4$ alors      $Tab=[-8,-4,1,2,5]$
\pause
\item $i=5$ alors     $Tab=[-8,-4,1,2,5,7]$ 
\end{itemize}
}

\end{frame}




\frame[containsverbatim]{
\frametitle{Insertion dans une liste triée en Python}


\begin{exampleblock}{}
\begin{verbatim}
def insertion(t,l):
    i=0
    taille=len(l)
    while (i<taille and t>l[i]):
        i+=1
    if (i<taille):    
        l[i+1:]=l[i:]
        l[i]=t
    else : 
        l.append(t)



insertion(5,l)
insertion(10,l)
print(l)
\end{verbatim}
\end{exampleblock}



}



\frame[containsverbatim]{
\frametitle{Suite de Fibonnaci}


La suite de Fibonacci est définie par la relation suivante :
\[
F_0 = 0, \quad F_1 = 1, \quad \text{et pour tout } n \ge 2, \quad F_n = F_{n-1} + F_{n-2}.
\]


On veut calculer toutes les éléments de la suite de Fibonnaci, jusqu'au $n$-ième. 

\vspace{0.5 cm}


\begin{alertblock}{Algorithme}
\begin{verbatim}
Input: n, valeurs initiales f0=0, f1=1
Output : Le tableau Fib contenant les nombres de Fibonnaci


Fib=[0,1]
i=2
Tant que  (i <=n) faire: 
      On rajoute Fib[i]=Fib[i-2]+Fib[i-1] dans le tableau
      i<-i+1
Renvoyer Fib

\end{verbatim}
\end{alertblock}




}


\frame[containsverbatim]{
\frametitle{Suite de Fibonacci}

On veut calculer le $n$-ième élément de la suite de Fibonacci. \\

On remarque qu'à chaque étape on a besoin de connaître que les deux valeurs précédentes. 

\begin{alertblock}{Algorithme}
\begin{verbatim}
Input: n, valeurs initiales a=0, b=1
Output : n-ième valeur dans la suite

i=2
Tant que  (i <= n) faire: 
      a,b<- b, a+b
      i<-i+1
Renvoyer a

\end{verbatim}
\end{alertblock}

}



\frame[containsverbatim]{
\frametitle{Suite de Fibonacci}



\begin{exampleblock}{En Python}
\begin{verbatim}
def fibonacci(n):
    if n < 0:
       print("n doit être >= 0")
    a, b = 0, 1
    for _ in range(n):
        a, b = b, a + b
    return a

\end{verbatim}
\end{exampleblock}


\vspace{0.5 cm}

\textbf{\`A retenir:} Un programme correct n'est n'est pas juste un programme qui fonctionne mais aussi: 
\begin{itemize}
\item qui utilise le moins de mémoire possible
\item qui résout le problème avec le moins d'opérations possibles!
\end{itemize}


}
\end{comment}





\end{document}
